% PAGES: 297
% Closes the sentence left open at the end of sec09_c3_3.tex (page 296 /
% folio 280 ends mid-word: "...and it would appear that VaR is"). No
% environment was open across that boundary, so this file simply resumes
% the same sentence, then finishes section 9.6.1 and opens section 9.7
% "References" (folio 281).
a subadditive measure; see an exercise at the end of this chapter. However,
formula (9.93) only holds if the portfolio return is normally distributed,
which is not the case in practice, and, in particular, is not the case for
the portfolios from the example above.

\section{References}

For an overview of the latest advances in asset allocation and portfolio
optimization see Fabozzi and Markowitz \cite{14}. Advanced risk and portfolio
management techniques are analyzed in depth in Meucci \cite{29}.

Efficient portfolios invested in only two assets are analyzed in detail in
many books; see. e.g., Ruppert \cite{33}.
% NOTE: source prints "see. e.g.," with an extra period after "see" (folio
% 281); reproduced as printed.

An expanded treatment of VaR can be found in Hull's risk management book
\cite{23}. An approach further tailored to practical applications can be
found in Alexander \cite{2}.

Expected shortfall was introduced as a measure of risk in the seminal paper
of Artzner et al. \cite{6}. Further details on expected shortfall (cVar) and
other alternatives to VaR are included in Andersen and Piterbarg \cite{4}.
% NOTE: the citation numbers printed on this page (Alexander \cite{2},
% Andersen and Piterbarg \cite{4}) differ from the numbers used for the same
% authors elsewhere in the already-transcribed book (e.g. ch06/ch07 cite
% Alexander as \cite{1} and Andersen and Piterbarg as \cite{3}). Both are
% reproduced exactly as printed on folio 281 -- not reconciled here -- and
% flagged for the bibliography cross-check pass.
